\subsection{Simplifying Ratios of Monomials} 
To simplify a rational expression in which the numerator and denominator are both monomials, we use the properties of exponents, specifically:
\begin{theorem}[Quotient of Powers Rule]
For any real number $x$ and positive integers $m\neq n$,
$$\mbox{If $m>n$, }\dfrac{x^m}{x^n}=x^{m-n}
\qquad
\mbox{If $m<n$, }\dfrac{x^m}{x^n}=\dfrac{1}{x^{n-m}}$$
\end{theorem}
We can apply this to each variable.
\begin{example}
Simplify.
\begin{lpic}{}{4}
\foreach \y/\l/\n/\d in {1/a/x^2y^3z/x^3yz^5,3/b/3a^2b^5c/9a^2c^4d^2} {
	\fractextright{1}{-\y}{\l. };
	\fractextleft{1}{-\y}{$\dfrac{\n}{\d}={}$};
}
\end{lpic}
\end{example}
While they aren't actually monomials, we can also do this with expressions that include negative exponents.
\begin{example}
Simplify. Give your answer without using negative exponents.
\begin{lpic}{}{2}
\fractextleft{1}{-1}{$\dfrac{x^{-3}y^2z^5}{x^2y^{-4}z^3}={}$};
\end{lpic}
\end{example}
Note that if a variable appears only once with a negative exponent, we can write it with a positive exponent by ``moving'' it between the numerator and denominator.
\begin{example}
Simplify. Give your answer without using negative exponents.
\begin{lpic}{}{2}
\fractextleft{1}{-1}{$\dfrac{4a^{-3}b}{5b^{-2}c^{-1}}={}$};
\end{lpic}
\end{example}
